\scanpage{41}
Or on a
\[
\varphi_f(x,v)=v(x)\qquad\deleted{\text{qui est}}
\]
d'où aussitôt la conclusion.

\uline{Faisceautisation} Pour tt ouvert $U$ de $X$, soit $\underline{\Phi}_X(U)$ l'ens des fonctions admissibles $\Gamma$-\unsure{invariantes} sur $U(\overline K)\times M_{\overline K}$.

Pour $U$ variable, on trouve ainsi un \uline{préfaisceau} \inserted{$\underline{\Phi}_X$} sur $X$, je dis que c'est un \uline{faisceau}. Évidemment, c'est un ``préfaisceau séparé'', reste à montrer que si une fonction
\[
\varphi:X(\overline K)\times M_{\overline K}\unsure{-|\longrightarrow}\mathbb R
\]
invariante par $\Gamma$ est telle que l'on a $\varphi|U_i(\overline K)\times M_{\overline K}$ admissible pour tt $i$ (où $(U_i)$ est un rec. ouvert de $X$) alors $\varphi$ est admissible. On est ramené au cas où $X$ est affine, et il semble qu'on ait besoin de savoir que \inserted{\unsure{déjà|alors|la réunion}} \uline{tte partie} $M_K$-\uline{bornée} \uline{de} $X(\overline K)\times M_{\overline K}$
% archive sheet number: 40
